\documentclass[11pt]{article}

\usepackage[a4paper,margin=2.35cm]{geometry}
\usepackage{booktabs}
\usepackage{array}
\usepackage{tabularx}
\usepackage{graphicx}
\usepackage{float}
\usepackage{caption}
\usepackage{xcolor}
\usepackage{enumitem}
\usepackage{amsmath}
\usepackage{hyperref}
\usepackage{microtype}

\definecolor{chapterblue}{HTML}{2F5F85}
\definecolor{textgray}{HTML}{333333}

\graphicspath{{../../charts/}}
\hypersetup{
  colorlinks=true,
  linkcolor=chapterblue,
  urlcolor=chapterblue,
  citecolor=chapterblue
}
\captionsetup{font=small,labelfont=bf}
\setlist[itemize]{leftmargin=1.6em,itemsep=0.15em}
\setlength{\parindent}{0pt}
\setlength{\parskip}{0.65em}
\renewcommand{\arraystretch}{1.18}

\title{\textbf{Part 5: Data Preparation, Descriptive Analysis, and Modeling Handoff}}
\author{DASE7337 Group Project: Airline Delay and Cancellation Operational Risk}
\date{April 2026}

\begin{document}
\maketitle

\begin{abstract}
This chapter follows the earlier discussion of JFK's aviation context, core delay factors, SWOT analysis, and PEST analysis. Its purpose is to explain how the BTS On-Time Performance data were transformed into modeling-ready operational risk assets. The emphasis is not only on what was cleaned, but also on why each cleaning decision matters: consistent field definitions, removal of non-operational records, separation of frequency and severity variables, and construction of descriptive outputs that support the later frequency, severity, state-dependent, and aggregate risk modeling stages.
\end{abstract}

\setcounter{section}{4}
\section{Data Preparation and Descriptive Analysis}

\subsection{Role of This Chapter}
The previous chapter shows that flight delay risk at JFK is shaped by both external shocks, such as weather and airspace constraints, and internal operational frictions, such as aircraft rotation, ground handling, and airline scheduling. The role of this chapter is to turn that qualitative risk language into a stable quantitative dataset. In that sense, data cleaning is not a background technical step. It is the bridge between managerial risk identification and the statistical models used later in the project.

The dataset used here is the BTS On-Time Performance Airline Delay Cause dataset. The sample is restricted to monthly airline-level arrival records at John F. Kennedy International Airport from January 2020 to December 2025. The raw dataset contains 570 airline-month records. After cleaning, 569 records remain. The final sample covers 72 airport months, 12 airlines, 657,671 arrival flights, 137,247 delayed arrivals of 15 minutes or more, 18,230 cancellations, 1,873 diversions, and 11,125,838 arrival delay minutes.

\subsection{Processing Workflow}
The processing workflow has four main steps. First, the raw CSV file and the official BTS column definition file are read and checked against the expected schema. Second, airline names and the airport name are standardized so that the same business entity is not split across slightly different labels. Third, raw BTS field names are rewritten into more readable analytical names. For example, \texttt{arr\_flights} becomes \texttt{total\_arrival\_flights}, and \texttt{arr\_del15} becomes \texttt{delayed\_arrivals\_15\_plus}. Fourth, the project exports a readable full dataset, a core cleaned dataset, a column dictionary, a monthly risk summary, a delay-cause summary, and an airline risk profile.

\begin{table}[H]
\centering
\caption{Cleaning and sample coverage}
\begin{tabular}{lr}
\toprule
Metric & Value \\
\midrule
Raw airline-month records & 570 \\
Cleaned airline-month records & 569 \\
Removed for missing or zero arrival flights & 1 \\
Removed for delayed arrivals exceeding total arrivals & 0 \\
Years covered & 2020--2025 \\
Airport months & 72 \\
Airlines & 12 \\
Total arrival flights & 657,671 \\
Total arrival delay minutes & 11,125,838 \\
\bottomrule
\end{tabular}
\end{table}

Two quality-control rules are especially important. The first rule removes records where total arrival flights are missing or equal to zero, because such records cannot produce meaningful delay rates, cancellation rates, or severity measures. The second rule checks whether delayed arrivals exceed total arrival flights, which would be a clear logical error. No such error appears in the current sample, but keeping the check makes the workflow safer when the dataset is updated later.

\subsection{Variable Design and Modeling Logic}
The cleaned core dataset preserves three groups of variables. The first group contains frequency variables, including total arrival flights, delayed arrivals, cancellations, and diversions. The second group contains severity variables, including total arrival delay minutes and cause-specific delay minutes for Airline, Weather, NAS, Security, and Late Aircraft. The third group contains identifiers, namely year, month, and airline name. This structure directly supports the later frequency--severity modeling framework.

At the airline-month level, the basic frequency measures are
\[
\mathrm{DelayRate}_{a,t} =
\frac{\mathrm{DelayedArrivals}_{a,t}}{\mathrm{TotalArrivals}_{a,t}},
\qquad
\mathrm{CancellationRate}_{a,t} =
\frac{\mathrm{CancelledArrivals}_{a,t}}{\mathrm{TotalArrivals}_{a,t}}.
\]
Severity is measured by average delay minutes per delayed arrival:
\[
\mathrm{AverageSeverity}_{a,t} =
\frac{\mathrm{TotalDelayMinutes}_{a,t}}{\mathrm{DelayedArrivals}_{a,t}}.
\]
If the denominator is zero, the airline-month should be treated as having no observed delay event rather than being forced into an artificial severity value. This separation is useful because operational risk modeling usually asks two different questions: how often does a disruption occur, and how severe is it once it occurs?

\subsubsection*{Delay-Equivalent Minutes as a Unified Loss Metric}
Because cancellations and diversions are not recorded in delay minutes by BTS, the project converts them into a common unit called \emph{delay-equivalent minutes}. The total operational impact for a given airline-month is defined as
\[
L_{a,t} = \mathrm{TotalDelayMinutes}_{a,t}
         + 180 \times \mathrm{CancelledArrivals}_{a,t}
         + 240 \times \mathrm{DivertedArrivals}_{a,t}.
\]
The choice of 180~minutes for cancellations and 240~minutes for diversions is grounded in regulatory thresholds, international passenger-rights conventions, and empirical evidence on recovery times.

\textbf{Why 180~minutes for cancellations.}
The U.S.\ Department of Transportation Tarmac Delay Rule (14~CFR~259.4), effective since April~29, 2010, prohibits airlines from keeping passengers on the tarmac for more than three hours (180~minutes) on domestic flights, with fines of up to USD~27,500 per passenger per violation \cite{dot2010tarmac}. This regulation creates a de facto boundary between a severe delay and a cancellation: as a flight approaches the three-hour mark, airlines frequently cancel rather than risk regulatory penalties. A DOT Office of Inspector General audit confirmed that cancellation rates rose significantly in the first three years after the rule took effect, validating the three-hour threshold as the practical tipping point between delay and cancellation \cite{dotoig2014tdr}. In parallel, EU~Regulation~261/2004 treats arrival delays exceeding three hours as equivalent to cancellations for the purpose of passenger compensation (EUR~250--600 depending on distance) \cite{eu261}, reinforcing the international consensus that three hours marks the boundary at which a delay becomes functionally indistinguishable from a cancellation. From a passenger-impact perspective, a cancelled flight at a major hub such as JFK typically requires rebooking on the next available service, with average wait times of two to four hours, making 180~minutes a conservative lower bound.

\textbf{Why 240~minutes for diversions.}
The same DOT tarmac rule sets a four-hour (240-minute) limit for international flights \cite{dot2010tarmac}. Because JFK is one of the busiest international gateways in the United States, a substantial share of its diversions involve international services, making the four-hour regulatory ceiling a natural reference point. More importantly, a diversion imposes strictly greater operational disruption than a cancellation: the aircraft has already departed, consumed fuel, and landed at an unplanned airport. The recovery process includes ground handling at the alternate airport, possible refueling, passenger re-screening, and either a continuation leg to the original destination or ground transportation. Peterson et~al.\ (2021) estimated that the average direct operating cost of a single diversion is approximately USD~38,000, with non-medical diversions taking significantly longer to recover than medical ones \cite{peterson2021diversion}. Four hours therefore represents a reasonable lower-bound estimate of the total additional delay imposed by a diversion. The ordering $240 > 180$ also preserves the intuitive severity hierarchy: a diversion is operationally more disruptive than a cancellation because it adds an airborne phase, extra fuel burn, crew-duty complications, and downstream schedule propagation that a pre-departure cancellation avoids.

In summary, the 180- and 240-minute penalties are anchored simultaneously to U.S.\ federal regulatory thresholds (DOT Tarmac Delay Rule), international passenger-rights practice (EU~261/2004), and empirical diversion-recovery research. This design gives the delay-equivalent-minutes metric both regulatory grounding and operational logic, while avoiding the introduction of subjective monetary assumptions.

\begin{thebibliography}{9}
\bibitem{dot2010tarmac}
U.S.\ Department of Transportation, ``Enhancing Airline Passenger Protections,'' 75~FR~32318, June~2010. 14~CFR~259.4.

\bibitem{dotoig2014tdr}
U.S.\ DOT Office of Inspector General, ``Effects of the Tarmac Delay Rule on Flight Cancellations and Delays,'' Report No.\ AV-2014-069, 2014.

\bibitem{eu261}
European Parliament and Council, Regulation~(EC) No~261/2004 establishing common rules on compensation and assistance to passengers in the event of denied boarding and of cancellation or long delay of flights, February~2004.

\bibitem{peterson2021diversion}
D.~C.~Peterson et~al., ``Data-Driven Estimation of the Impact of Diversions Due to In-Flight Medical Emergencies on Flight Delay and Aircraft Operating Costs,'' \emph{Aerospace Medicine and Human Performance}, vol.~92, no.~2, pp.~81--87, 2021.
\end{thebibliography}

\subsection{Descriptive Evidence}
Figure \ref{fig:monthly-trend-en} shows the monthly trend in delayed and cancelled arrivals from 2020 to 2025. The highest delay month is July 2023, with 3,657 delayed arrivals and a delay rate of 32.69\%. The highest cancellation month is March 2020, with 1,900 cancellations, reflecting the exceptional shock at the beginning of the pandemic. This contrast is important: delays and cancellations are related disruptions, but they do not always move together. Delays often reflect congestion and operational pressure, while cancellations may also reflect large-scale planning changes and external shocks.

\begin{figure}[H]
\centering
\includegraphics[width=0.94\linewidth]{chart_1_monthly_trend.png}
\caption{Monthly delay and cancellation trend at JFK}
\label{fig:monthly-trend-en}
\end{figure}

The delay-cause structure is highly concentrated. Airline and Late Aircraft account for 36.16\% and 35.14\% of total delay minutes, respectively. NAS contributes 24.32\%, while Weather and Security together account for less than 5\%. This finding connects directly to the risk identification section: weather is a visible aviation risk, but in this sample the main day-to-day risk mass comes from the internal operating chain and system congestion.

\begin{figure}[H]
\centering
\includegraphics[width=0.92\linewidth]{chart_2_delay_cause_breakdown.png}
\caption{Delay-minute contribution by cause}
\label{fig:cause-breakdown-en}
\end{figure}

\begin{table}[H]
\centering
\caption{Main delay-cause shares}
\begin{tabular}{lrr}
\toprule
Cause & Delay minutes & Share \\
\midrule
Airline & 4,022,961 & 36.16\% \\
Late Aircraft & 3,910,112 & 35.14\% \\
NAS & 2,705,624 & 24.32\% \\
Weather & 466,088 & 4.19\% \\
Security & 21,053 & 0.19\% \\
\bottomrule
\end{tabular}
\end{table}

Figure \ref{fig:airline-profile-en} summarizes airline-level risk profiles by combining delay rate, cancellation rate, and average delay severity. JetBlue Airways contributes the largest total delay minutes, Hawaiian Airlines has the highest delay rate, and Envoy Air has the highest average delay minutes per delayed flight. This should not be read as a simple ranking of airline performance, because airlines differ in sample size, route structure, and operating role at JFK. It is more useful as a modeling guide, showing which airline-month observations may contribute more strongly to tail outcomes.

\begin{figure}[H]
\centering
\includegraphics[width=0.88\linewidth]{chart_3_airline_risk_profile.png}
\caption{Airline-level risk profile}
\label{fig:airline-profile-en}
\end{figure}

\subsection{How This Supports Later Modeling}
This chapter supports the later quantitative analysis in three ways. First, the core cleaned dataset separates frequency variables from severity variables, allowing the next stage to compare candidate frequency distributions such as Poisson and Negative Binomial, and candidate severity distributions such as Lognormal and Weibull. Second, the cause summary helps translate the BTS database taxonomy into the operational risk language used in the course: Airline and Late Aircraft can be grouped as internal operating-chain risk, NAS as system constraint risk, and Weather plus Security as external shock risk. Third, the monthly summary preserves chronological order, which makes it possible to define normal and disrupted operating states and then estimate state transitions.

The main contribution of this chapter is therefore not a single chart. It is the construction of an interpretable, reproducible, and model-ready risk dataset. It connects the managerial analysis in the earlier chapter with the statistical modeling that follows, so the final models are not just mathematical exercises but structured operational risk analysis.

\subsection{Limitations and Next Step}
Two limitations remain. First, the data are monthly airline-level records, so individual flights, hourly airport congestion, and specific weather events cannot be directly observed. Short-term disruption propagation is therefore captured only indirectly through monthly aggregates. Second, BTS allocates some cause-specific delay counts according to reporting rules, so these fields should be read as attribution measures rather than perfectly independent event counts.

Even with these limits, the prepared dataset is strong enough for the quantitative requirements of the group project. The next chapter can build directly on these outputs to compare frequency distributions, compare severity distributions, introduce state dependence, and simulate annual aggregate operational risk under different scenarios.

\end{document}
